\specialchapter{符号索引}

 \(\mathcal{C}_{\mathrm{B}}(\mathrm{X};\mathrm{X}^{\prime})\) (B-态射) \ldots.. 2 \(\mathrm{Isom}_{\mathrm{B}}(\mathrm{X};\mathrm{X}^{\prime})\) (B-同构) \ldots.. 2 \(\mathrm{Aut}_{\mathrm{B}}(\mathrm{X})\) (B-自同构) .. 2 \(\mathrm{X}_{b}\) (纤维) \ldots.. 2 \(\mathrm{X}_{\mathrm{A}}\) (诱导空间) \ldots.. 2 \(\mathrm{X} \times_{\mathrm{B}} \mathrm{X}^{\prime}\) (纤维积) \ldots.. 3 \(\Delta_{\mathrm{X}}\) (对角线) \ldots.. 4 \(\delta_{\mathrm{X}}\) (对角态射) \ldots.. 4 \(f^{*}(\mathrm{X})\) (换基) .. 5 \(\prod_{\mathrm{B}} \mathrm{X}_{i}\) (纤维积) \ldots.. 5 \(\mathcal{C}_{c}(\mathrm{U};\mathrm{V})\) (从 U 到 V 的连续映射) \ldots.. 19 \(\mathcal{S}(\mathrm{A};\mathrm{E})\) (连续截面) .. 33 \(s \mid \mathrm{U}\) (限制) \ldots.. 43 \(\underline{\mathcal{F}}(\mathrm{B};\mathrm{X})\) (映射层) \ldots.. 45 \(\underline{\mathcal{C}}(\mathrm{B};\mathrm{X})\) (连续映射层) \ldots.. 45 \(\underline{\mathcal{L}}(\mathrm{B};\mathrm{E}), \underline{\mathcal{L}}(\mathrm{E})\) (截面层) \ldots.. 45 \(\underline{\mathcal{C}}_{\mathrm{B}}(\mathrm{E};\mathrm{E}^{\prime})\) (B-态射层) \ldots.. 45 \(\underline{\mathcal{Mor}}_{\mathrm{B}}(\mathrm{E};\mathrm{E}^{\prime})\) (B-态射层) \ldots.. 46 \(\underline{\mathcal{Isom}}_{\mathrm{B}}(\mathrm{E};\mathrm{E}^{\prime})\) (同构层) \ldots.. 46 \(\underline{\mathcal{C}}^{r}(\mathrm{X};\mathrm{Y})\) (\(C^{r}\) 类映射的层) \ldots.. 46 \(\underline{\mathfrak{P}}(\mathrm{B})\) (子空间层) \ldots.. 46 \(\mathrm{E}_{\mathcal{F}}\) (平展空间) \ldots.. 49 \(\widetilde{\mathcal{F}}\) (预层的伴随层) 53 \(u_{*}(\mathcal{F})\) (预层、层的直像) 57 \(u^{*}(\mathcal{G})\) (预层、层的逆像) 58 \(\mathcal{G}_{\mathrm{A}}\) (诱导层) 59 \(\mathrm{G}^{\circ}\) (反群) 91 \(\mathcal{C}_{\mathrm{B}}^{\mathrm{G}}(\mathrm{E};\mathrm{E}^{\prime})\) (主纤维丛的态射) 92 \(\mathrm{E} \times^{\mathrm{G}} \mathrm{F}\) (相伴纤维丛) . 106 \(\mathrm{Z}_{\mathrm{cont}}^{1}(\mathrm{B},\mathcal{U},\mathrm{G})\) (连续 1-上闭链) 114 \(\mathrm{H}_{\mathrm{cont}}^{1}(\mathrm{B},\mathcal{U},\mathrm{G})\) (上同调类) 116 \(\mathrm{P}(\mathrm{B},\mathrm{G})\) (主丛的同构类) 118 \(\mathrm{P}(\mathrm{B},\mathcal{U},\mathrm{G})\) (主丛的同构类) 118 \(\mathrm{H}_{\mathrm{cont}}^{1}(\mathrm{B},\mathrm{G})\) (上同调类) 119 \(\mathcal{C}((\mathrm{X},x);(\mathrm{Y},y))\) (带基点连续映射) 120 \(p_{\mathrm{A}}\) (凸部分 A 的规度) 123 \(\mathrm{S}(\mathrm{X})\) (空间 X 的悬垂) 149\\
Som(C) (箭图的顶点) 152\\
Fl(C) (箭图的箭) . 152

 \(o_{C}\) (箭图的源映射) \ldots.. 152 \(t_{C}\) (箭图的靶映射) \ldots.. 152 \(\mathrm{Fl}_{a,b}(\mathrm{C}), \mathrm{C}_{a,b}\) (连接 a 与 b 的箭) \ldots.. 152 Som(φ) (箭图态射) \ldots.. 152 \(\mathrm{Fl}(\varphi)\) (箭图态射) \ldots.. 152 \(c*d\) (拼接道路) \ldots.. 154 \(\pi_{0}(\mathrm{C})\) (连通分支) \ldots.. 155 \(\pi_{0}(\varphi)\) (在连通分支上诱导的映射) .. 155 \(\overline{f}\) (Fl(C) 上的对合) \ldots.. 155 Ch(C) (箭图中的道路集) \ldots.. 160 \(\Omega_{C}\) (箭图的道路范畴) \ldots.. 160 \(f^{-1}\) (可逆箭的逆) \ldots.. 161 \(G_{a}\) (迷向群) \ldots.. 162 Int(f) (迷向群的同构) \ldots.. 162 \(\varphi_{a}\) (迷向群的同态) \ldots.. 162 Orb(G) (群胚的轨道) \ldots.. 162 Orb(φ) (通过取轨道由 φ 诱导的映射) . 162 \(\mathcal{B}(X,p)\) (置换群胚) \ldots.. 163 \(\varinjlim_{i}G_{i}\) (群胚的归纳极限) \ldots.. 164 X ×s X (映射的配对群胚) \ldots.. 165 \(\varphi^{*}G\) (逆像群胚) \ldots.. 166 Ker(φ) (群胚态射的核) \ldots.. 166 G/H (商群胚) \ldots{} 170 \(\overline{\varphi}\) (通过取商诱导的群胚态射) . 171 \(\varpi_{G}\) (道路类群胚) \ldots.. 172 \(\varpi(\varphi)\) (道路类群胚态射) \ldots.. 173 F(A) (自由群) \ldots.. 173 Grp(C) (自由群胚) \ldots. 174

G/F (由箭缩并得到的群胚) \ldots{} 176 π₁(G, a) (图在 a 处的庞加莱群) \ldots{} 178 u ∼ v (群胚态射的同伦) \ldots{} 181 Coh(φ, ψ) (上同伦子) \ldots{} 184 * Gᵢ (群的自由积) i∈I 193 Coeg(φ, ψ) (余均衡子) \ldots{} 199 C(X; Y) (从 X 到 Y 的连续映射) \ldots{} 230 {[}X; Y{]} (连续映射的同伦类) . 230 f (连续映射的同伦类) \ldots{} 230 C((X, x); (Y, y)) (带基点连续映射) \ldots{} 232 Sₙ (欧氏球面) \ldots{} 234 Cyl(f) (映射的柱) \ldots{} 237 X ∪ f B (附着所得拓扑空间) \ldots{} 248 sₓ/A (X/A 的基点) \ldots{} 252 Côn(f) (映射的锥) 253 c * d (拼接道路) \ldots{} 256 ¯ c (逆道路) \ldots{} 256 Λ(X) (道路空间) \ldots{} 257 Λₓ(X) (起点为 x 的道路) . 257 λₓ,y(X) (起点为 x、终点为 y 的道路) \ldots{} 257 π₀(X) (道路连通分支) \ldots{} 259 π₀(f) (在道路连通分支上由 f 诱导的映射) \ldots{} 259 ϖₓ,y(X) (从 x 到 y 的道路类) \ldots{} 290 Ω(X) (环路空间) \ldots{} 290 Ωₓ(X) (x 处的圈) \ldots{} 290 εₓ (常值圈) \ldots{} 290 ϖ(X) (庞加莱群胚) 292 π₁(X, x) (基本群) 292 ϖ(f) (f 诱导的庞加莱群胚态射) \ldots{} 294 f*(γ) (ϖₓ,y(f)(γ) 的缩写) \ldots{} 294

 \(x \cdot \gamma\) (庞加莱群胚的作用) \ldots.. 304 \(e_{\mathrm{B}}\) (靶映射) \ldots.. 342 \(\varepsilon_{\mathrm{B}}\) (靶映射) \ldots.. 342 \(\pi_{1}(\mathrm{G})\) (拓扑群的庞加莱群) \ldots.. 375 \(\mathrm{R}_{\tau} \{\xi z_{1}, z_{2}\}\) (与下降数据相关联的等价关系) \ldots.. 383

\((Z/R_{\tau},q)\) (由下降数据商得的空间 \((Z,p)\)) \ldots. 384 \(C_{\tau,\tau'}(Z;Z')\) (与下降数据相容的态射) \ldots. 387 \textbar G\textbar{} (图的几何实现) \ldots. 417 \(\bigvee_{i}(X_{i},x_{i})\) (带基点拓扑空间的楔和) \ldots{} 421

